% Appendix E — N6 ternary X2MH6 2/2 falsification; cation-VEC insufficient.
\section{Appendix E --- The $X_2MH_6$ ternary funnel: cation-VEC is
necessary but insufficient}
\label{app:E}

The N5 wall (Appendix~\ref{app:D}) motivated a ternary funnel: stuffed-cation
octahedral hydrides $X_2MH_6$ (Fm$\bar{3}$m, K$_2$PtCl$_6$ prototype) that
might decouple coupling from stability. This appendix reports the
\emph{pre-registered} falsification of that route at ambient pressure.

\paragraph{The cation-VEC rule.}
The valence-electron count of the octahedral unit is the closed-form
\begin{equation}
\mathrm{VEC} = 2\,v(X) + v(M) + 6,
\qquad
n_{\sigma^\ast} = \mathrm{VEC} - 18,
\end{equation}
where $v(\cdot)$ is the group valence and $n_{\sigma^\ast}$ is the filling of
the \ce{H6}-derived $\sigma^\ast$ antibonding manifold. The model predicts a
sweet spot at VEC$\in\{19,20\}$ ($n_{\sigma^\ast}=1$--2), where $N(E_F)$ is
maximised. Both \ce{Mg2IrH6} ($2\cdot2+9+6=19$) and \ce{Li2CuH6}
($2\cdot1+11+6=19$) land exactly on it --- so the rule, as a heuristic, is
self-consistent. The question is whether VEC$=19$ is \emph{sufficient}.

\paragraph{Pre-registered falsifiers.}
We registered the pass/fail threshold before harvesting either result:
a candidate is dynamically stable (and the sweet spot ``sufficient'') only
if every $q$-point clears $\texttt{min\_freq} > -\SI{50}{\per\centi\meter}$
with at most one hard-imaginary mode. Tab.~\ref{tab:ternary} is the verdict.

\begin{table}[h]
\centering
\small
\begin{tabular}{lcccl}
\toprule
\textbf{Candidate} & \textbf{VEC} & \textbf{ambient stable?} &
  \textbf{min freq (\si{\per\centi\meter})} & \textbf{verdict} \\
\midrule
\ce{Mg2IrH6} & 19 & \textbf{no} (48\% hard-imag) & $-2235$ & \tierRed\ FALSIFIED \\
\ce{Li2CuH6} & 19 & \textbf{no} ($8$ hard-imag modes) & $-944.92$ & \tierRed\ FALSIFIED \\
\bottomrule
\end{tabular}
\caption{The $X_2MH_6$ cation-VEC$=19$ sweet spot, falsified 2/2 at ambient
pressure. \ce{Mg2IrH6}: 5/13 partial $q$-verdict already violates the
threshold (PR \#247). \ce{Li2CuH6}: the second $q$-point alone locks 8
hard-imaginary modes (PR \#275). The K$_2$PtCl$_6$ Fm$\bar{3}$m prototype is
ambient-unstable for both.}
\label{tab:ternary}
\end{table}

\paragraph{The falsifier statements (verbatim, F-N6 ledger).}
\begin{quote}\small
\textbf{F-N6-1.} If \ce{Mg2IrH6} Fm$\bar{3}$m at $P=0$ satisfies
$\texttt{min\_freq} > -\SI{50}{\per\centi\meter}$ on all $q$ with
hard-imag count $\le 1$, the VEC$=19$ sweet spot is \emph{sufficient};
otherwise it is only \emph{necessary}.\\
\textbf{Result:} \tierRed\ PASSED 2026-05-26 --- threshold clearly violated
($-\SI{2235}{\per\centi\meter}$, 48\% hard-imag). Sufficiency falsified.

\medskip
\textbf{F-N6-2.} The same threshold for \ce{Li2CuH6}; failure
family-falsifies the K$_2$PtCl$_6$ prototype at ambient.\\
\textbf{Result:} \tierRed\ PASSED 2026-05-27 --- $q_2$ alone: 8 hard-imag,
$-\SI{944.92}{\per\centi\meter}$. $X_2MH_6$ Fm$\bar{3}$m 2/2 falsified.
\end{quote}

\paragraph{Reading.}
The cation-VEC rule survives as a necessary filter (it correctly places both
candidates on the sweet spot and explains why \ce{Mg}$\to$\ce{Ca}/\ce{Sr}
substitution kills $T_c$ despite preserving VEC --- the larger ionic radius
expands the lattice and softens the H phonons). But VEC$=19$ does \emph{not}
guarantee dynamical stability: the octahedral $\sigma^\ast$ filling that
maximises $N(E_F)$ also drives the soft modes that collapse the lattice. This
is the N5 wall reappearing in the ternary family. The honest consequence is
that the ambient $X_2MH_6$ octahedral route, as tested, is closed; the
clathrate $MXH_8$ route (\ce{LaBeH8} anchor) remains open and is the
campaign's recommended next probe (\S\ref{sec:conclusion}).
